% ---------------------------------------------------------------------
\section{Proposed Method}\label{sec:method}
% ---------------------------------------------------------------------

As seen in \Fig{fig:architecture}, \method{} has three parts that act in sequence at every step.
An attention policy maps a set of charger tokens to one raw action per charger, and a safety layer turns the raw actions into commanded powers that respect per-EV bounds and the transformer limit.
During training, a policy-gradient method updates the policy from the realised profit with the safety layer treated as part of the environment, and the same layer acts unchanged during deployment.

\begin{figure}[t]
\centering
\figorbox{figures/fig_architecture.pdf}{\linewidth}{6cm}
\caption{Architecture and control loop of \method{}. Solid arrows form the control loop used in training and deployment; dashed arrows are used in training only. The safety layer (shaded) is an exact projection that keeps the transformer limit and keeps each EV's requirement reachable while the EV receives its lower bound (Theorem~\ref{thm:invariance}).}
\label{fig:architecture}
\end{figure}

% ---------------------------------------------------------------------
\subsection{Token representation and attention policy}\label{sec:policy}

The observation is a set of $N$ charger tokens $\xi_i(t)\in\R^{12}$ and one global token $\xi_0(t)\in\R^{38}$, and no feature encodes the index or the position of a charger.
A charger token holds occupancy, energy and capacity data, the remaining required energy and steps, the laxity of \Sec{sec:bounds}, the power limits and the bounds $l_i$ and $u_i$.
The global token holds prices for the next $H_f=8$ steps, time features, the current limits, load and PV forecasts and four station ratios.
All features are normalised so that they do not grow with $N$, and Table~\ref{S-tab:features} lists them.

Each token is embedded linearly with one of two learned type embeddings, $e_i=W_c\xi_i+b_c+v_c$ for chargers and $e_0=W_g\xi_0+b_g+v_g$ for the global token.
The sequence $[e_0;e_1;\dots;e_N]$ passes through $n_L=3$ pre-LayerNorm transformer encoder layers \cite{vaswani2017attention,ba2016layer,xiong2020layer} of width $d_m=128$ with $n_H=4$ heads and no positional encoding.
A key-padding mask removes empty chargers as attention keys, and the global key is never masked.
Then a shared actor head maps each charger output $h_i$ to the mean $\mu_{\theta,i}$ and standard deviation $\sigma_{\theta,i}$ of a Gaussian raw action, where $\theta$ denotes the actor parameters.
The policy is $\pi_\theta(a\mid s)=\prod_{i\in\mathcal N(t)}\mathrm{Normal}\big(a_i;\mu_{\theta,i}(s),\sigma_{\theta,i}^2(s)\big)$.
A separate critic encoder with parameters $\omega$ pools all unmasked token outputs with one learned attention query \cite{lee2019set} into the value $V_\omega(s)$.

The construction is permutation-equivariant, masked tokens do not influence the outputs of unmasked tokens, and the number of parameters does not depend on $N$ (Lemma~\ref{S-lem:equivariance}).
Whatever the number of chargers, the actor has 610,370 parameters and the critic 676,481.
Weights trained at $N\in\{10,20\}$ can therefore run at $N=200$, but performance there is not guaranteed, because attention over 200 tokens averages over more keys than seen in training.

% ---------------------------------------------------------------------
\subsection{Per-EV bounds from reachable energy}\label{sec:bounds}

For each connected EV, the safety layer needs the least power that keeps its departure requirement reachable, and we drop the index $i$ where no confusion arises.
Full-power charging for one step maps the energy $E$ to $\varphi(E)=E+\Delta t\,\eta^{\sm c}\,\bar P(E)$.
Assumption A1 of \Sec{sec:invariance} asks that $\varphi$ is continuous and non-decreasing, so that starting a step with more energy never leaves less energy after full-power charging.
The backward-reachable thresholds, the least energies from which the requirement is still reachable at full power, are
\begin{equation}\label{eq:thresholds}
B(0)=\Ereq,\qquad B(k)=\min\big\{E\in[E^{\min},E^{\sm{cap}}]:\ \varphi(E)\ge B(k-1)\big\},\quad k\ge1 .
\end{equation}
Under A1, an EV is individually feasible at step $t$, in the sense that some admissible control meets \Eq{eq:requirement}, exactly when $E(t)\ge B(\tau(t))$ (Lemma~\ref{S-lem:reach}).
Without A1 this characterisation fails, because charging less now could leave more reachable energy later.
For EV2Gym's exponential taper, Lemma~\ref{S-lem:taper} proves A1 for every step length.

For a connected EV with energy $E=E_i(t)$ and $\tau=\tau_i(t)$ remaining actions, the bounds of the safety layer are
\begin{equation}\label{eq:bounds}
u_i(t)=\bar P(E),\qquad
l_i(t)=\begin{cases}
\min\Big\{u_i(t),\ \max\Big\{\pmn(E),\ g^{-1}\Big(\dfrac{\Bh(\tau-1)-E}{\Delta t}\Big)\Big\}\Big\}, & E\ge\Bh(\tau),\\[2mm]
\bar P(E), & E<\Bh(\tau),
\end{cases}
\end{equation}
where $\Bh$ is a stored threshold table and the outer minimum is a numerical guard.
In the first case, $l_i$ is the smallest admissible power after which $\tau-1$ full-power steps still reach $\Ereq$, and an individually infeasible EV is charged at maximum power.
A negative lower bound allows V2G discharging while slack remains, and as seen in \Fig{fig:concepts}(a), the bound keeps the energy on or above the latest-start trajectory.

\begin{figure}[t]
\centering
\figchain{\linewidth}{\analysisdir/figures/fig2_concepts.pdf}{\analysisdir/figures_conceptual/fig_concepts.pdf}{figures/fig2_concepts.pdf}
\caption{Concepts of the safety layer. (a) Energy of one EV during its stay, with the latest-start trajectory $\Bh_i$ and the per-step options $\Delta t\,u_i$, $\Delta t\,l_i$ and $\Delta t\,\pmn_i$. (b) Projection for two EVs onto the box $[l,u]$ cut by the slab $-\Crev\le p_1+p_2\le\Ctil$.}
\label{fig:concepts}
\end{figure}

The table $\Bh_i(k)$, $k=0,\dots,\tau_i(t_i^{\sm{arr}})$, is computed once at arrival by bisection on $\varphi_i$, so the bounds cost $O(1)$ per EV and step.
Each stored table is certified in the sense that
\begin{equation}\label{eq:certified}
\Bh(0)\ge\Ereq,\qquad \varphi\big(\Bh(k)\big)\ge\Bh(k-1),\qquad \Bh(k)\le\max\{\Bh(k-1),E^{\min}\},\qquad \Bh(k)\in[E^{\min},E^{\sm{cap}}],
\end{equation}
and the proofs use only these properties.
The laxity $\chi_i(t)=\big(E_i(t)-\Bh_i(\tau_i(t))\big)/\big(\eta_i^{\sm c}\Pnom_i\Delta t\big)$, measured in steps, is non-negative exactly when the EV is individually feasible.

% ---------------------------------------------------------------------
\subsection{Projection}\label{sec:projection}

Let $x\in\R^N$ be the raw action in kW, with $x_i=l_i=u_i=0$ for empty chargers.
The safety layer uses the tightened forward limit $\Ctil(t)=C(t)-n(t)\varepsilon_q$, where the margin $\varepsilon_q=0.008$ kW per charger absorbs the actuation rounding of the simulator (Section~\ref{S-app:actuation}).
Its feasible set is $\mathcal K(t)=\{p\in\R^N:\ l\le p\le u,\ -\Crev(t)\le\mathbf 1^\top p\le\Ctil(t)\}$, and the safety layer is the Euclidean projection
\begin{equation}\label{eq:proj}
\Pi_t(x)=\arg\min_{p\in\mathcal K(t)}\tfrac12\|p-x\|_2^2 .
\end{equation}

\begin{proposition}[Projection]\label{prop:projection}
\begin{enumerate}
\item[(i)] The set $\mathcal K(t)$ is non-empty if and only if $l\le u$, $\sum_il_i\le\Ctil$, $\sum_iu_i\ge-\Crev$ and $-\Crev\le\Ctil$.
\item[(ii)] If $\mathcal K(t)\neq\emptyset$, the unique minimiser of \Eq{eq:proj} is $p_i^{*}=\clip(x_i-\lambda^{*},l_i,u_i)$. The common shift $\lambda^{*}$ is zero if the clipped point already satisfies the aggregate limit, and otherwise it is the root of $\lambda\mapsto\sum_i\clip(x_i-\lambda,l_i,u_i)$ at the violated limit.
\end{enumerate}
\end{proposition}

The proof verifies the Karush--Kuhn--Tucker (KKT) conditions, and the supplementary material also gives the Jacobian (Proposition~\ref{S-prop:jacobian}).
As seen in \Fig{fig:concepts}(b), the projection shifts all free coordinates by the same amount and clips them until the aggregate limit holds.
One sort of the $2N$ breakpoints $\{x_i-u_i\}\cup\{x_i-l_i\}$ and one pass over them give $\lambda^{*}$, so the layer costs $O(N\log N)$ per step.
A post-check with a bisection fallback keeps the aggregate limit under floating-point rounding, and Algorithm~\ref{S-alg:safety} lists one complete step.

The algebra is not new: \Eq{eq:proj} is the continuous quadratic knapsack problem \cite{kiwiel2007breakpoint,wang2022efficient}, and the sLLF rate rule uses a closely related form, a common level scaled per EV by its peak rate \cite{chen2022smoothed}.
Our use differs in its input, a learned per-charger action of a size-agnostic policy, and in the lower bound $l_i$ from backward reachability with taper and V2G sign range.
Moreover, the aggregate constraint is two-sided, and the layer acts during training.

% ---------------------------------------------------------------------
\subsection{Feasibility preservation}\label{sec:invariance}

Intuitively, the safety layer never lets an EV fall below the least energy from which its requirement is still reachable at full power, and the common shift keeps the station inside the transformer band.
Only the first property can fail, namely when the transformer cannot give every EV its lower bound at once.
We call a step feasible when $\mathcal K(t)$ is non-empty, and on an infeasible step a fallback rule replaces the projection.
The conditions of the result are collected below, and only the first five are assumptions.
Technically, the induction behind it extends the laxity argument of Chen et al. \cite{chen2022smoothed} to the taper, V2G discharging, a two-sided coupling and arbitrary raw actions.

\begin{center}
\textbf{Conditions of Theorem~\ref{thm:invariance}}\\[1mm]
\footnotesize
\setlength{\tabcolsep}{3pt}
\begin{tabular}{@{}l>{\raggedright}p{6.4cm}>{\raggedright\arraybackslash}p{6.2cm}@{}}
\toprule
Label & Condition & Status in this study \\
\midrule
A1 & $\varphi$ continuous and non-decreasing & proved for EV2Gym's taper (Lemma~\ref{S-lem:taper}) \\
A2 & departure step $d_i$ and requirement $\Ereq_i$ known at arrival & relaxed by the conformal bound (\Sec{sec:conformal}) \\
A3 & one-step model: $E_i(t+1)\ge E_i(t)+\Delta t\,g_i(p_i)$ and $p_i\le\tilde p_i\le p_i+\varepsilon_q$ for commands within \Eq{eq:limits} & proved for EV2Gym's actuation (Lemma~\ref{S-lem:a3}); may fail on hardware \\
A4 & feasible arrival: $E_i^{\sm{arr}}\ge\Bh_i(\tau_i(t_i^{\sm{arr}}))$ & enforced by the request pipeline \\
A6 & grid headroom: $C(t)\ge N\varepsilon_q$ and $\Crev(t)\ge0$ & exogenous; checked through the fallback counts \\
\midrule
A5 & feasible step: $\sum_il_i\le\Ctil$, $\sum_iu_i\ge-\Crev$, $\Ctil\ge-\Crev$ & a closed-loop event, not assumed by (a) and (b); measured \\
\bottomrule
\end{tabular}
\end{center}

A6 uses the number of chargers $N$ instead of $n(t)$, so that it does not depend on the policy, and it fails in our scenarios with a probability below $2\times10^{-9}$ per episode (Section~\ref{S-app:a6}).

\begin{theorem}[Feasibility preservation]\label{thm:invariance}
Let the commanded powers be $p(t)=\Pi_t(x(t))$ on feasible steps and the output of the fallback rule otherwise, for an arbitrary sequence of finite raw actions, during training or deployment.
\begin{enumerate}
\item[(a)] Under A3 and A6, the realised powers satisfy the transformer limit $-\Crev(t)\le\sum_i\tilde p_i(t)\le C(t)$, and the commanded powers satisfy the physical limits of \Eq{eq:limits}, at every step.
\item[(b)] Under A1--A4, every EV $i$ with $p_i(t)\ge l_i(t)$ throughout its stay meets $E_i(d_i+1)\ge\Ereq_i$, in particular every EV whose stay contains no infeasible step.
\item[(c)] If, in addition, every step is feasible, every EV meets its departure requirement.
\end{enumerate}
\end{theorem}

\begin{proof}[Proof sketch]
On feasible steps, Proposition~\ref{prop:projection}(i) gives $p(t)\in\mathcal K(t)$, and A3 transfers the tightened band to the realised powers.
Under A6, the fallback rule can only meet the case of excess required power described below, and its output keeps the commanded sum at $\Ctil(t)$.
For (b), a one-step lemma (Lemma~\ref{S-lem:onestep}) carries the invariant $E_i(t)\ge\Bh_i(\tau_i(t))$ from arrival, where A4 gives it, through every step with $p_i\ge l_i$, and at the last step it gives $E_i(d_i+1)\ge\Bh_i(0)\ge\Ereq_i$.
Part (c) follows from (b), and Section~\ref{S-app:theorem} gives the full proof.
\end{proof}

The transformer part depends neither on the policy nor on A1 or A5, whereas the departure part is a per-EV statement.
A5 can fail even when a joint plan exists.
Two EVs that each need 10 kWh within two 1-h steps under a 10 kW limit can share the capacity equally.
A policy that idles in the first step, however, leaves $\sum_il_i=20$ kW in the second step.
We therefore measure how often EVs meet infeasible steps.

% ---------------------------------------------------------------------
\subsection{Fallback rule}\label{sec:fallback}

On an infeasible step, the fallback rule replaces the projection with a deterministic least-laxity-first allocation.
In the relevant case, the required power $\sum_il_i$ exceeds $\Ctil$.
Every EV then starts at its physical lower limit $\pmn_i$, and the remaining budget up to $\Ctil$ goes to the EVs in increasing order of laxity, each up to $l_i$.
This output keeps the transformer limit, and every feasible EV that still receives $l_i$ stays feasible (Lemma~\ref{S-lem:fallback}).
Three further cases cover contradictory limits, excess PV export and a negative budget, and A6 excludes all three (Section~\ref{S-app:fallback}).
For the EVs that receive less than $l_i$, the total departure shortfall of an episode is bounded a priori by the accumulated deficit of these steps (Lemma~\ref{S-lem:deficit}).
Admission control in the request pipeline makes infeasible steps rarer, but it relies on forecasts and does not enter Theorem~\ref{thm:invariance}.

% ---------------------------------------------------------------------
\subsection{Conformal lower bound on the departure step}\label{sec:conformal}

To relax A2 for the departure step, we use split conformal prediction (CP) \cite{papadopoulos2002inductive,vovk2022algorithmic,lei2018distributionfree,angelopoulos2023conformal} in a one-sided, quantile-regression form \cite{romano2019conformalized}.
A gradient-boosting regressor $\hat q$ estimates a lower quantile of $d_i$ from the arrival step, the weekday, the arrival energy, the requested energy and the EV model.
Calibration draws one session uniformly from each of $n_{\sm{cal}}$ independently seeded episodes and uses the scores $e_j=\hat q(\mathbf w_j)-d_j$.
With $\hat Q_{1-\alpha}$ the $\lceil(n_{\sm{cal}}+1)(1-\alpha)\rceil$-th smallest score, the bound $\hat d_i^{\sm{lo}}=\max\{t_i^{\sm{arr}},\lceil\hat q(\mathbf w_i)-\hat Q_{1-\alpha}\rceil\}$ replaces $d_i$ in the bounds, the request pipeline and the features, and after $\hat d_i^{\sm{lo}}$ the layer uses one remaining step.

\begin{proposition}[Conformal departure guarantee]\label{prop:conformal}
Let the calibration sessions and the test session be exchangeable.
Then $\mathbb P(d_i\ge\hat d_i^{\sm{lo}})\ge1-\alpha$, and under A1, A3 and A4 with $\hat d_i^{\sm{lo}}$ in place of $d_i$, $\mathbb P\big(E_i(d_i+1)\ge\Ereq_i\big)\ge1-\alpha-\mathbb P\big(\exists\,t\le d_i:\ p_i(t)<l_i(t)\big)$.
\end{proposition}

Here $\Ereq_i$ is the requirement re-certified with $\hat d_i^{\sm{lo}}$, and the probability is over the calibration data and a uniformly chosen EV of a test episode, so the guarantee is marginal.
Calibration needs data from each test configuration, so the conformal variant is not zero-shot.

% ---------------------------------------------------------------------
\subsection{Training algorithm}\label{sec:training}

The policy is trained with proximal policy optimization (PPO) \cite{schulman2017proximal} and generalised advantage estimation \cite{schulman2016highdimensional}, with the safety layer, the fallback rule and the actuation treated as part of the environment.
PPO therefore uses the likelihood ratio of the raw Gaussian sample.
The vanilla policy-gradient estimator, equal to the PPO surrogate gradient at the old parameters, is then unbiased for the return of the projected actions.
For stochastic policies, projection inside the environment and inside the policy are equivalent in this sense \cite{markgraf2025safe}, and the choice avoids the zero gradients of an active projection \cite{kasaura2023benchmarking}.
Tokens with $u_i-l_i\le10^{-9}$ kW cannot influence any commanded power and are dropped from the log-probability and the entropy.
Fallback steps are masked out of the actor loss but kept for the critic.
Many raw actions map to the same projected point, so on feasible steps a regulariser with weight $\beta=0.01$ keeps the policy mean close to $\mathcal K(t)$,
\begin{equation}\label{eq:reg}
\mathcal L_{\sm{reg}}(\theta)=\beta\,\mathbb E_{s}\Big[\frac1{n(t)}\sum_{i\in\mathcal N(t)}\Big(\frac{\bar x_i-\Pi_t(\bar x)_i}{P^{\sm{ref}}}\Big)^{2}\Big],\qquad \bar x=P^{\sm{ref}}\mu_\theta(s).
\end{equation}
By the Moreau identity, the gradient of $\tfrac12\mathrm{dist}^2(x,\mathcal K)$ is $x-\Pi_t(x)$, so $\Pi_t(\bar x)$ is evaluated under stop-gradient \cite{chen2021enforcing,markgraf2025safe}.
Each rollout runs eight environments with $N=10$ and eight with $N=20$, and Algorithm~\ref{S-alg:ppo} gives the full training loop.
